%% ELEC2013 2007

%% ----------------------------------------------------------------
%% exam.tex
%% ----------------------------------------------------------------
\documentclass[answers]{sotonExamBoxes}    % use option [answers] to include them
%\documentclass[answers]{uosexamb}
%% ----------------------------------------------------------------
\usepackage{graphicx}
%\usepackage{pstricks}
\usepackage{alltt}
\usepackage{bm}
\usepackage{amssymb}
\usepackage{amsfonts}
\usepackage{amsmath}
\thinmuskip=1mu
%\renewcommand{\ttdefault}{pcr}
\renewcommand{\ttdefault}{pcr}
\newcommand{\tr}{\textsf{T}}
\newcommand{\e}[1]{{\rm e}^{#1}}
\newcommand{\E}[1]{{\rm e}^{{\displaystyle #1}}}
\newcommand{\logg}[1]{\log\left(#1\right)}
\newcommand{\Prob}[1]{\mathbb{P}\left(#1\right)}
\newcommand{\Step}{\mathop{\mathrm{O}\hspace{-15pt}\rule[8.5pt]{11pt}{1.2pt}\hspace{2pt}}}
\newcommand\diag{\mathop{\mathrm{diag}}}
\newcommand{\grad}{\bm{\nabla}}
\newcommand{\dd}{\mathrm{d}}
\newcommand{\class}{\mathcal{C}}
\newcommand{\data}{\mathcal{D}}
\DeclareMathAlphabet{\mat}{OT1}{cmss}{bx}{n}
\newcommand{\pd}[2]{\ensuremath{\frac{\qpartial #1}{\qpartial #2}}\xspace}
\newcommand{\M}[1]{\ensuremath{\boldsymbol{#1}}\xspace}
\newcommand{\bst}{\rule{0mm}{5mm}}
\newcommand{\av}[2][\,]{\mathbb{E}_{#1}\left[ {\strut #2} \right]}
\newcommand{\len}[1]{\| #1 \|}
\graphicspath{{figures/} {../2018-19/figures}}
\newcommand{\inner}[2]{\left\langle #1, #2\right\rangle}
\graphicspath{{figures/}}
\newcommand{\squeeze}{\setlength{\itemsep}{-2pt}}
\newcommand{\bra}[1]{\left( #1 \right)}
\usepackage{parskip}
\AddToHook{env/samepage/begin}{%
\AddToHook{para/before}{\nopagebreak}%
\AddToHook{para/after}{\nopagebreak}%
}
% \usepackage{picture}
\begin{document}
\unitTitle{Vector Space Problem Sheet}
\authors{Adam Pr\"ugel-Bennett}
\unitcode{COMP6208}
\semester{2}
\year{2022}
\duration{45}{0}


\maketitle


% ***************************** question 1 **************************


\begin{question}{20}
  \begin{qparts}
    
      \qpart[3]{Show that if the columns, $\bm{m}_j$, of an $n\times n$ matrix,
        $\mat{M}$, are not independent then $\mat{M}$ does not have a
        unique inverse.  For the columns to not be independent implies
        there exists some set of coefficients $c_j$ such that
        \begin{align*}
          \sum_{i=1}^n c_j\, \bm{m}_j = \bm{0}.
        \end{align*}
        We note that the $i^{th}$ element of $\bm{m}_j = M_{ij}$ so
        \begin{align*}
          \sum_{i=1}^n c_j\, M_{ij} = 0
        \end{align*} 
      \begin{answer}
        We consider the vector $\bm{c}$ with elements $c_i$, then the
        $i^{th}$ component of $\mat{M}\,\bm{c}$ is equal to
        \begin{align*}
          [\mat{M}\,\bm{c}]_i = \sum_{i=0}^n M_{ij}\, c_j = 0
        \end{align*}
        or $\mat{M}\, \bm{c} = \bm{0}$.  Consequently for any vector
        $\bm{x}$ and any scalar $b$ we have that
        \begin{align*}
          \mat{M}\, \bra{\strut \bm{x} + b\, \bm{c}} =  \mat{M}\, \bm{x}
        \end{align*}
        if we denote $\bm{y} = \mat{M}\, \bm{x}$, then there is not a
        unique vector $\bm{x}$ such that $\bm{y} = \mat{M}\, \bm{x}$,
        but rather a family of vectors $\bm{x} + b\, \bm{c}$ for any $b$.
      \end{answer}}{\lines{5}}
    
        
      \qpart[3]{Show that for three $n\times n$ matrices $\mat{A}$,
        $\mat{B}$ and $\mat{C}$ the inverse of
        $\mat{A}\,\mat{B}\,\mat{C}$ is
        $\mat{C}^{-1}\mat{B}^{-1}\mat{A}^{-1}$.  You can do this by
        multiplying the product of matrices, but its inverse and use
        associativity of matrix multiplication
      \begin{answer}
        \begin{align*}
          \bra{\mat{C}^{-1}\mat{B}^{-1}\mat{A}^{-1}} \bra{\strut
          \mat{A}\,\mat{B}\,\mat{C}}
          &=
            \mat{C}^{-1}\mat{B}^{-1}\bra{\mat{A}^{-1}\mat{A}}\,\mat{B}\,\mat{C}
          \\
          &= \mat{C}^{-1}\bra{\mat{B}^{-1}\mat{B}}\,\mat{C} =
            \mat{C}^{-1}\mat{C} = \mat{I}
        \end{align*}
      \end{answer}}{\lines{3}}
    

    
      \qpart[3]{For non-square matrices, $\mat{A}$ and $\mat{B}$ we can
        have $\mat{A}\,\mat{B} = \mat{I}$ while $\mat{B}\,\mat{A} \neq
        \mat{I}$.  Show this is the case for the matrices
        \begin{align*}
          \mat{A} &=
                    \begin{pmatrix}
                      1 & 0
                    \end{pmatrix}
          & \mat{B} &=
                      \begin{pmatrix}
                        1 \\ 0
                      \end{pmatrix}
        \end{align*}
      \begin{answer}
        \begin{align*}
          \mat{A}\,\mat{B} &=
                             \begin{pmatrix}
                               1 & 0
                             \end{pmatrix}
                             \begin{pmatrix}
                               1 \\ 0
                             \end{pmatrix}
                             = (1) = \mat{I}_1
          &
            \mat{B}\,\mat{A} &=
                               \begin{pmatrix}
                                 1 \\ 0
                               \end{pmatrix}
                               \begin{pmatrix}
                                 1 & 0
                               \end{pmatrix} 
                               =
                               \begin{pmatrix}
                                 1 & 0 \\ 0 & 0
                               \end{pmatrix}
        \end{align*}
      \end{answer}}{\lines{5}}

    
      \qpart[2]{For finite-dimensional square matrices if
        $\mat{A}\mat{B}= \mat{I}$ then $\mat{B}\mat{A}= \mat{I}$. Thus
        $\mat{B} = \mat{A}^{-1}$.  This is not entirely trivial to
        prove.  We break the proof down.

        Show that if $\mat{A}\mat{B}= \mat{I}$ then $\mat{B}$ is
        injective (one-to-one).  This mean that for two different
        vectors $\bm{x}$ and $\bm{y}$ then
        $\mat{B}\, \bm{x} \neq \mat{B} \bm{y}$.  Equivalently
        $\mat{B}\, (\bm{x} - \bm{y}) \neq \bm{0}$, unless $\bm{x}=\bm{y}$,
        but as $\bm{x}$ and $\bm{y}$ are arbitrary the only vector for
        which $\mat{B}\,\bm{x}=\bm{0}$ is the vector $\bm{x}=\bm{0}$.
        Prove that if $\mat{A}\mat{B}= \mat{I}$ and
        $\mat{B}\,\bm{x}=\bm{0}$ then $\bm{x}=\bm{0}$.
      \begin{answer}
        Suppose $\mat{B}\,\bm{x}=\bm{0}$ then multiply by $\mat{A}$
        \begin{align*}
          \mat{A}\, \mat{B}\,\bm{x} &= \bm{0} \\
          \mat{I}\,\bm{x} &= \bm{0}
        \end{align*}
        Thus $\bm{x} = \bm{0}$.
      \end{answer}}{\lines{8}}
    

    
      \qpart[5]{Next we want to show that in a finite dimensional space
        a square matrix that is injective is also surjective (onto).
        That is, if $\mat{B}$ is injective then for any vector $\bm{y}$
        there exists a vector $\bm{x}$ such that
        $\bm{y}= \mat{B}\bm{x}$.  The reason a finite dimensional square
        matrix that is injective is also surjective is not so obvious.
        To prove this we consider a set of $n$ \textbf{basis} vectors
        that span the entire space $\mathbb{R}^n$.  We call these
        $\bm{v}_1$, $\bm{v}_2$, \ldots, $\bm{v}_n$.  Because these are
        independent the only set of numbers $c_i$ such that
        $\sum_i c_i \bm{v}_i=\bm{0}$ are when $c_1=c_2=\ldots=c_n=0$.
        Show that the set of vectors $\bm{v}_i' = \mat{B}\, \bm{v}_i$
        form a set of basis vectors (i.e. they are linear
        independent).  This is easily proved by contradiction.
      \begin{answer}
        Since $\mat{B}$ is injective and $\bm{v}_i\neq \bm{0}$ (since
        basis vectors are non-zero), clearly $\bm{v}_i' = \mat{B}\,
        \bm{v}_i$ are non-zero.  The non-trivial part is to show these
        vectors are linear independent.  Suppose this is no the case and
        there exists a set of coefficients $c_i'$ such that
        \begin{align*}
          \sum_{i=1}^n c_i'\, \bm{v}_i' &= \bm{0} \\
          \sum_{i=1}^n c_i'\, \mat{B} \, \bm{v}_i &= \bm{0} \\
          \mat{B} \, \sum_{i=1}^n c_i'\, \bm{v}_i &= \bm{0}.      
        \end{align*}
        But, as $\mat{B}$ is injective this implies
        \begin{align*}
          \sum_{i=1}^n c_i'\, \bm{v}_i &= \bm{0}.
        \end{align*}
        Since the vectors $\bm{v}_i$ are linear independent, so
        this is a contradiction.
      \end{answer}}{\lines{8}}
    

    
      \qpart[2]{Use the fact that the vectors $\bm{v}_i'$ are linearly
        independent to show that, for any vector $\bm{y}$, we can find a
        vector $\bm{x}$ such that $\bm{y} = \mat{B}\,\bm{x}$
      \begin{answer}
        For any vector $\bm{y}$ can be represented by a weighted sum of
        $n$ independent vectors.  In particular we can write
        \begin{align*}
          \bm{y} &= \sum_{i=1}^n c_i\, \bm{v}'_i = \sum_{i=1}^n c_i\,
                   \mat{B} \, \bm{v}_i  \\
                 &= \mat{B} \bra{\sum_{i=1}^n c_i\, \bm{v}_i} = \mat{B}\,\bm{x}
        \end{align*}
        where $x = \sum_{i=1}^n c_i\, \bm{v}_i$.
      \end{answer}}{\lines{6}}
    

    
      \qpart[2]{Using the fact that $\mat{A}\,\mat{B}= \mat{I}$ and
        $\mat{B}$ is surjective show $\mat{B}\,\mat{A} =
        \mat{I}$
      \begin{answer}
        We not that for all $\bm{x}$
        \begin{align*}
          \mat{A}\,\mat{B}\,\bm{x} &= \bm{x} \\
          \mat{B}\,\mat{A}\,\mat{B}\,\bm{x} &= \mat{B}\,\bm{x} 
        \end{align*}
        But as $\mat{B}$ is surjective for any vector $\bm{y}$ there
        exists a vector $\bm{x}$ such that $\bm{y} = \mat{B}\,\bm{x}$.
        Thus,
        \begin{align*}
          \mat{B}\,\mat{A}\,\bm{y} = \bm{y},
        \end{align*}
        but as this is true for eany vector $\bm{y}$ it follows that $\mat{B}\,\mat{A}=\mat{I}$.
      \end{answer}}{\lines{5}}
    
  \end{qparts}
\end{question}
\freshpage

% ***************************** question 2 **************************

\begin{question}{20}
  We want to compute the derivative of $\logg{\det(\mat{X})}$ with
  respect to the matrix $\mat{X}$.  Our strategy will be to use the
  fact that
  \begin{align*}
    \lim_{\epsilon\rightarrow 0} \frac{\logg{\det(\mat{X} + \epsilon\, \mat{U})} -
    \logg{\det\bra{\mat{X}}}}{\epsilon} = \mathrm{tr} \, \mat{U}^\tr \frac{\partial
    \logg{\det(\mat{X})}}{\partial \mat{X}}.
  \end{align*}
  We therefore expand $\logg{\det(\mat{X}+ \epsilon\, \mat{U}))}$ to
  first order in $\epsilon$, subtract $\logg{\det\bra{\mat{X}}}$, divide by
  $\epsilon$, rearrange and read off $\frac{\partial
    \logg{\det(\mat{X})}}{\partial \mat{X}}$.  However, rearranging in
  this case is quite complicated.
  
  \begin{qparts}
    \qpart[5]{Using $\det(\mat{A}\,\mat{B}) =
      \det(\mat{A})\,\det(\mat{B})$ and $\log(a\,b)= \log(a) +
      \log(b)$ expand $\logg{\det(\mat{X}+ \epsilon \mat{U}))}$ as a
      sum of $\logg{\det(\mat{X})}$ and another logarithm.
      \begin{answer}
        \begin{align*}
          \logg{\det\bra{\strut \mat{X}+\epsilon\,\mat{U}}}
          &= \logg{\det\bra{\mat{X}\,\bra{\mat{I} + \epsilon\,\mat{X}^{-1}\mat{U}}}}\\
          &= \logg{\det(\mat{X}) \, \det(\mat{I} + \epsilon\,\mat{X}^{-1}\mat{U})}\\
          &= \logg{\det(\mat{X})} + \logg{\det\bra{\mat{I} +
            \epsilon\,\mat{X}^{-1}\mat{U}}}
        \end{align*}
      \end{answer}}{\lines{5}}

    \qpart[2]{The hard part is to expand $\det\bra{\mat{I} +
        \epsilon\,\mat{M}}$. Write out $\det\bra{\mat{I} +
        \epsilon\,\mat{M}}$ explicitly for a 2 
      dimensional matrix and expand to first order in $\epsilon$
      \begin{answer}
        \begin{align*}
          \det\bra{\mat{I}+\epsilon\,\mat{M}}
          &=
            \det\begin{pmatrix}
              1+\epsilon\, M_{11} & \epsilon\, M_{12}\\
              \epsilon\, M_{21} & 1+\epsilon\, M_{22}
            \end{pmatrix}
            = (1+\epsilon\, M_{11})\,(1+\epsilon\, M_{22}) -
            \epsilon^2\, M_{21}\,M_{12}\\
          &= 1 + \epsilon \bra{M_{11} + M_{22}} + O(\epsilon^2)
        \end{align*}
      \end{answer}}{\lines{8}}

    \qpart[5]{For an $n\times n$ matrix $\mat{M}$ compute the
      determinant of
      \begin{align*}
        \det\bra{\mat{I} + \epsilon\, \mat{M}} = \det
        \begin{pmatrix}
          1 + \epsilon\, M_{11} & \epsilon\, M_{12} & \ldots & \epsilon\,M_{1n}  \\
          \epsilon \, M_{21} & 1 + \epsilon\, M_{22} & \ldots &\epsilon\,M_{2n}\\
          \vdots & \vdots & \ddots & \vdots \\
          \epsilon \, M_{n,1} &  \epsilon\, M_{n2} & \ldots &1+\epsilon\,M_{nn}
        \end{pmatrix}
      \end{align*}
      up to first order in $\epsilon$.  Recall we can compute the
      determinant a matrix $\mat{X}$ as
      \begin{align*}
        \det(\mat{X}) = \sum_{i=1}^n (-1)^j\, X_{1j} \, \det(\hat{\mat{X}}_{1j})
      \end{align*}
      where $\hat{\mat{X}}_{1j}$ is the matrix obtained
      by removing row 1 and column $j$ from the matrix $\mat{X}$.
      \begin{answer}
        Expanding the determinant we get a term equal to the
        product of the diagonal elements
        \begin{align*}
          \prod_{i=1}^n (1+\epsilon \, M_{ii}) = 1 +\epsilon\,
          \sum_{i=1}^n M_{ii} + O(\epsilon^2)
        \end{align*}
        all other terms in the determinant involve at least terms of
        order $\epsilon^2$ so do not contribute to order $\epsilon$.
        In consequence
        \begin{align*}
          \det\bra{\mat{I} + \epsilon\, \mat{M}} = 1 + \epsilon
          \mathrm{tr} \, \mat{M} + O(\epsilon^2)
        \end{align*}
      \end{answer}}{\lines{8}}

    \qpart[3]{For a more elegant proof we can use the fact that
      \begin{align*}
        \det(\mat{M}) &= \prod_{i=1}^n \lambda_i& \mathrm{tr}\,\mat{M} &=
                                                                \sum_{i=1}^n \lambda_i
      \end{align*}
      where $\lambda_i$ are the eigenvalues of the matrix $\mat{M}$.
      We will prove these results later when we look at eigen-systems.
      Note that if $\bm{v}_i$ is an eigenvector of the matrix
      $\mat{M}$, with eigenvalue $\lambda_i$ then it is an eigenvector
      of $\mat{I} + c\, \mat{M}$ with eigenvalue $1+c\,\lambda_i$
      since
      \begin{align*}
        (\mat{I} + c\, \mat{M})\, \bm{v}_i = \bm{v}_i + c\,\lambda_i\,
        \bm{v}_i
        = (1+c\,\lambda_i)\, \bm{v}_i.
      \end{align*}
      Using these results show that $\det\bra{\mat{I} + \epsilon\, \mat{M}} = 1 + \epsilon\,
      \mathrm{tr} \, \mat{M} + O(\epsilon^2)$.
      \begin{answer}
        \begin{align*}
          \det\bra{\mat{I} + \epsilon\, \mat{M}}
          &= \prod_{i=1}^n (1+\epsilon\, \lambda_i)
            = 1 + \epsilon\, \sum_{i=1}^n \lambda_i + O(\epsilon^2) \\
          &= 1 + \epsilon\, \mathrm{tr}\, \mat{M} + O(\epsilon^2) \\
        \end{align*}
      \end{answer}}{\lines{5}}

    \qpart[5]{Using the results
      \begin{align*}
        \logg{\det(\mat{X} + \epsilon\, \mat{U})} - \logg{\det\bra{\mat{X}}}
        &= \logg{\det\bra{\mat{I} + \epsilon\,\mat{X}^{-1}\mat{U}}} \\
        \det\bra{\mat{I} + \epsilon\, \mat{M}} &= 1 + \epsilon
      \mathrm{tr} \, \mat{M} + O(\epsilon^2)
      \end{align*}
      together with the fact that $\log(1+x) = x + O(x^2)$, 
      compute $\frac{\partial \logg{\det(\mat{X})}}{\partial
        \mat{X}}$.  You may use the fact that $\log(1+x) = x +
      O(x^2)$ and $\mathrm{tr}\, \mat{A} = \mathrm{tr}\, \mat{A}^\tr$.
      \begin{answer}
        We note that
        \begin{align*}
          \logg{\det(\mat{X} + \epsilon\, \mat{U})} - \logg{\det\bra{\mat{X}}}
          &= \logg{1 + \epsilon\, \textrm{tr} \, \mat{X}^{-1}\mat{U} +
          O(\epsilon^2)}
          = \epsilon\,  \textrm{tr} \, \mat{X}^{-1}\mat{U} +
          O(\epsilon^2) = \epsilon\, \textrm{tr} \mat{U}^\tr (\mat{X}^{-1})^\tr
        \end{align*}
        Thus,
        \begin{align*}
          \frac{\partial \logg{\det(\mat{X})}}{\partial \mat{X}} =  (\mat{X}^{-1})^\tr
        \end{align*}
      \end{answer}}{\lines{7}}
  \end{qparts}
\end{question}

% ***************************** question 3 **************************



% ***************************** End of exam *****************************
\end{document}




